\documentclass[10pt,a4paper]{article}
% ============================================================
%  COURS : LA FONCTION INVERSE — 4 VOLETS — TOUT EN PORTRAIT
%  Compilation : pdflatex x2
% ============================================================
\usepackage[utf8]{inputenc}
\usepackage[T1]{fontenc}
\usepackage[french]{babel}
\usepackage{lmodern}
\usepackage{amsmath,amssymb,amsthm}
\usepackage[margin=0.6cm]{geometry}
\usepackage[table]{xcolor}
\usepackage{booktabs,array,tabularx,multirow}
\usepackage{enumitem}
\usepackage{multicol}
\usepackage{tikz}
\usetikzlibrary{positioning,arrows.meta,shapes.geometric,shapes.misc,
	calc,backgrounds,decorations.pathreplacing,fit}
\usepackage{hyperref}
\hypersetup{colorlinks=true,linkcolor=blue!70!black,urlcolor=blue,linktoc=all}
\usepackage{pdfcomment}
\usepackage[most]{tcolorbox}
\usepackage{pifont}

% ---- Silence des warnings de bookmarks ----
\pdfstringdefDisableCommands{%
	\def\frac#1#2{#1/#2}\def\textbf#1{#1}\def\textit#1{#1}%
	\def\quad{ }\def\qquad{ }\def\\{ }\def\,{ }\def\:{}%
	\def\to{->}\def\mapsto{->}\def\infty{inf}\def\pm{+/-}%
	\def\mathbb#1{#1}\def\lim{lim}\def\cup{U}\def\not{}%
	\def\searrow{dec}\def\nearrow{cro}\def\cdot{*}\def\le{<=}\def\ge{>=}%
	\def\sqrt#1{sqrt(#1)}\def\left{}\def\right{}%
}

% ---- Macros interactives ----
\newcommand{\pop}[2]{\pdfmarkupcomment[markup=Highlight,color=yellow!40]{#1}{#2}}
\newcommand{\tip}[2]{\pdftooltip{#1}{#2}}

% ---- Palette ----
\definecolor{cParCoeur}{RGB}{180,30,30}
\definecolor{cParCoeurBg}{RGB}{255,240,240}
\definecolor{cComprendre}{RGB}{30,80,160}
\definecolor{cComprendreBg}{RGB}{235,242,252}
\definecolor{cMethode}{RGB}{30,130,60}
\definecolor{cMethodeBg}{RGB}{235,250,240}
\definecolor{cPiege}{RGB}{220,110,20}
\definecolor{cPiegeBg}{RGB}{255,245,230}
\definecolor{cRappel}{RGB}{120,50,160}
\definecolor{cRappelBg}{RGB}{245,238,252}
\definecolor{cExemple}{RGB}{60,60,60}
\definecolor{cExempleBg}{RGB}{248,248,248}

% ---- Style commun des boîtes ----
\tcbset{
	boxrule=1pt, arc=5pt,
	left=7pt, right=7pt, top=5pt, bottom=5pt,
	fonttitle=\bfseries\sffamily\small,
	coltitle=white,
	enhanced, breakable,
}

\newtcolorbox{parcoeur}[1][]{
	colback=cParCoeurBg, colframe=cParCoeur,
	title={\ding{72}~\textbf{PAR C\OE UR}~\textbullet~\`a m\'emoriser}, #1}
\newtcolorbox{acomprendre}[1][]{
	colback=cComprendreBg, colframe=cComprendre,
	title={\ding{110}~\textbf{\`A COMPRENDRE}~\textbullet~sans m\'emoriser}, #1}
\newtcolorbox{methode}[1][]{
	colback=cMethodeBg, colframe=cMethode,
	title={\ding{43}~\textbf{M\'ETHODE}~\textbullet~\`a savoir appliquer}, #1}
\newtcolorbox{piege}[1][]{
	colback=cPiegeBg, colframe=cPiege,
	title={\ding{53}~\textbf{PI\`EGE}~\textbullet~\`a \'eviter absolument}, #1}
\newtcolorbox{rappel}[1][]{
	colback=cRappelBg, colframe=cRappel,
	title={\ding{45}~\textbf{RAPPEL}~\textbullet~niveau ant\'erieur}, #1}
\newtcolorbox{exemplebox}[1][]{
	colback=cExempleBg, colframe=cExemple,
	title={\ding{52}~\textbf{EXEMPLE}}, #1}

% ---- Titres de volets ----
\newcommand{\volet}[2]{%
	\begin{tcolorbox}[colback=blue!12, colframe=blue!60!black,
		fonttitle=\Large\bfseries\sffamily, coltitle=white,
		title={#1}]
		\centering\itshape #2
\end{tcolorbox}}

\title{\textbf{\color{blue!70!black}La fonction inverse}\\[4pt]
	\large Dérivée, variations, asymptotes, prix unitaire\\
	et combinaison avec les polynômes\\[6pt]
	\normalsize Cours complet en 4 volets}
\author{}\date{}

\begin{document}
	\maketitle
	
	% ============================================================
	% LÉGENDE
	% ============================================================
	\begin{tcolorbox}[colback=white, colframe=blue!50!black,
		title={\Large\bfseries\sffamily Comment lire ce cours ?}]
		\begin{center}
			\begin{tabular}{p{0.22\linewidth}p{0.7\linewidth}}
				\textcolor{cParCoeur}{\rule{10pt}{10pt}}~\textbf{PAR C\OE UR} & Formules, définitions, résultats \`a \emph{mémoriser mot pour mot}. \\[4pt]
				\textcolor{cComprendre}{\rule{10pt}{10pt}}~\textbf{\`A COMPRENDRE} & Idées, interprétations, raisonnements \`a \emph{comprendre sans r\'eciter}. \\[4pt]
				\textcolor{cMethode}{\rule{10pt}{10pt}}~\textbf{M\'ETHODE} & Proc\'edures, techniques \`a \emph{savoir appliquer} pas \`a pas. \\[4pt]
				\textcolor{cPiege}{\rule{10pt}{10pt}}~\textbf{PI\`EGE} & Erreurs fr\'equentes \`a \emph{\'eviter absolument}. \\[4pt]
				\textcolor{cRappel}{\rule{10pt}{10pt}}~\textbf{RAPPEL} & Pr\'erequis des ann\'ees ant\'erieures (\emph{\`a revoir si besoin}). \\[4pt]
				\textcolor{cExemple}{\rule{10pt}{10pt}}~\textbf{EXEMPLE} & Illustrations concr\`etes du cours.
			\end{tabular}
		\end{center}
		\vspace{4pt}
		\textbf{Interactif.} Les termes \pop{surlignés en jaune}{Ceci est une annotation interactive. Cliquez pour ouvrir la modale ; vous pouvez la d\'eplacer (glisser-d\'eplacer).} sont \textbf{cliquables}. Dans les volets 3 et 4, \emph{survolez} les n\oe uds pour afficher leur info-bulle.
	\end{tcolorbox}
	
	\tableofcontents
	\newpage
	
	% ============================================================
	% VOLET 1
	% ============================================================
	\volet{Volet 1 — Présentation classique}{Le cours sous sa forme traditionnelle : définitions, théorèmes, démonstrations et exemples.}
	
	\section{Rappels de collège et de seconde}
	
	\begin{rappel}[title={\ding{45}~\textbf{RAPPEL}~\textbullet~D\'efinition d'une fonction}]
		Une \textbf{fonction} $f$ associe \`a chaque r\'eel $x$ d'un ensemble $D$ (appel\'e \emph{ensemble de d\'efinition}) un unique r\'eel $f(x)$. On note $f:D\to\mathbb{R}$, $x\mapsto f(x)$.
	\end{rappel}
	
	\begin{rappel}[title={\ding{45}~\textbf{RAPPEL}~\textbullet~Sens de variation}]
		Soit $f$ d\'efinie sur un intervalle $I$.
		\begin{itemize}[leftmargin=*,itemsep=1pt]
			\item $f$ est \textbf{croissante} sur $I$ si $a<b \Rightarrow f(a)\le f(b)$.
			\item $f$ est \textbf{d\'ecroissante} sur $I$ si $a<b \Rightarrow f(a)\ge f(b)$.
			\item \guillemotleft~Strictement~\guillemotright{} : remplacer $\le$ (resp.\ $\ge$) par $<$ (resp.\ $>$).
		\end{itemize}
	\end{rappel}
	
	\begin{acomprendre}[title={\ding{110}~\textbf{\`A COMPRENDRE}~\textbullet~Lien d\'eriv\'ee / variations}]
		Si $f$ est d\'erivable sur $I$ : $f'>0 \Rightarrow f$ strictement croissante ; $f'<0 \Rightarrow f$ strictement d\'ecroissante ; $f'=0 \Rightarrow f$ constante. La r\'eciproque (variation $\Rightarrow$ signe de $f'$) est vraie aux points o\`u $f'$ existe.
	\end{acomprendre}
	
	\begin{rappel}[title={\ding{45}~\textbf{RAPPEL}~\textbullet~Asymptotes}]
		La droite $x=a$ est \textbf{asymptote verticale} si $\lim_{x\to a} f(x)=\pm\infty$.\\
		La droite $y=b$ est \textbf{asymptote horizontale} si $\lim_{x\to\pm\infty} f(x)=b$.\\
		La droite $y=mx+p$ est \textbf{asymptote oblique} si $\lim_{x\to\pm\infty}[f(x)-(mx+p)]=0$.
	\end{rappel}
	
	\section{D\'efinition de la fonction inverse}
	
	\begin{parcoeur}[title={\ding{72}~\textbf{PAR C\OE UR}~\textbullet~D\'efinition de la fonction inverse}]
		On appelle \pop{fonction inverse}{La fonction inverse associe a tout reel non nul son inverse multiplicatif. On note f de x egale 1 sur x.} la fonction
		\[
		f:\mathbb{R}^*\longrightarrow\mathbb{R},\qquad x\longmapsto f(x)=\frac{1}{x}.
		\]
		Son ensemble de d\'efinition est $\mathbb{R}^*=\mathbb{R}\setminus\{0\}$.
	\end{parcoeur}
	
	\begin{parcoeur}[title={\ding{72}~\textbf{PAR C\OE UR}~\textbullet~Parit\'e}]
		Pour tout $x\neq 0$, $f(-x)=-f(x)$ : la fonction inverse est \textbf{impaire}. Sa courbe (l'hyperbole) est sym\'etrique par rapport \`a l'origine $O$.
	\end{parcoeur}
	
	\section{D\'eriv\'ee}
	
	\begin{parcoeur}[title={\ding{72}~\textbf{PAR C\OE UR}~\textbullet~D\'eriv\'ee de la fonction inverse}]
		La fonction inverse est d\'erivable sur $\mathbb{R}^*$ et
		\[
		\boxed{\;f'(x)=-\frac{1}{x^{2}}\quad \text{pour tout } x\neq 0.\;}
		\]
	\end{parcoeur}
	
	\begin{methode}[title={\ding{43}~\textbf{M\'ETHODE}~\textbullet~Deux d\'emonstrations}]
		\textbf{M\'ethode 1 (r\`egle de puissance).} $f(x)=x^{-1}$ donne $f'(x)=-1\cdot x^{-2}=-1/x^{2}$.\\[4pt]
		\textbf{M\'ethode 2 (taux d'accroissement).} Pour $h\neq 0$,
		\[
		\frac{f(x+h)-f(x)}{h}=\frac{1}{h}\Bigl(\frac{1}{x+h}-\frac{1}{x}\Bigr)
		=\frac{-1}{x(x+h)}\xrightarrow[h\to 0]{}-\frac{1}{x^{2}}.
		\]
	\end{methode}
	
	\begin{parcoeur}[title={\ding{72}~\textbf{PAR C\OE UR}~\textbullet~G\'en\'eralisation : d\'eriv\'ee de $1/v$}]
		Si $v$ est d\'erivable sur $I$ et ne s'annule pas sur $I$, alors
		$\left(\dfrac{1}{v}\right)'=-\dfrac{v'}{v^{2}}$.
	\end{parcoeur}
	
	\section{Tableau de variations}
	
	\begin{parcoeur}[title={\ding{72}~\textbf{PAR C\OE UR}~\textbullet~Variations de la fonction inverse}]
		Puisque $f'(x)=-1/x^{2}<0$ sur $\mathbb{R}^*$, $f$ est \textbf{strictement d\'ecroissante sur chacun} des intervalles $\,]-\infty,0[\,$ et $\,]0,+\infty[$.
		
		\begin{center}\renewcommand{\arraystretch}{1.4}
			\begin{tabular}{c|ccccc}
				\toprule
				$x$ & $-\infty$ & & $0$ & & $+\infty$\\
				\midrule
				$f'(x)$ & & $-$ & $\|$ & $-$ & \\
				\midrule
				$f$ & $0$ & $\searrow$ & $\|$ & $\searrow$ & $0$\\
				\bottomrule
			\end{tabular}
		\end{center}
	\end{parcoeur}
	
	\begin{piege}[title={\ding{53}~\textbf{PI\`EGE}~\textbullet~D\'ecroissance \guillemotleft~globale~\guillemotright}]
		$f$ n'est \textbf{pas} d\'ecroissante sur $\mathbb{R}^*$ tout entier. En effet $-1<1$ mais $f(-1)=-1<f(1)=1$, ce qui contredit la d\'ecroissance globale (on aurait d\^u avoir $f(-1)\ge f(1)$).
	\end{piege}
	
	\section{Limites et asymptotes}
	
	\begin{parcoeur}[title={\ding{72}~\textbf{PAR C\OE UR}~\textbullet~Limites}]
		\[
		\lim_{x\to+\infty}\frac{1}{x}=0,\quad
		\lim_{x\to-\infty}\frac{1}{x}=0,\quad
		\lim_{x\to 0^{+}}\frac{1}{x}=+\infty,\quad
		\lim_{x\to 0^{-}}\frac{1}{x}=-\infty.
		\]
	\end{parcoeur}
	
	\begin{parcoeur}[title={\ding{72}~\textbf{PAR C\OE UR}~\textbullet~Asymptotes de l'hyperbole}]
		La courbe de la fonction inverse admet :
		\begin{itemize}[leftmargin=*,itemsep=1pt]
			\item une \pop{asymptote verticale}{L'axe des ordonnees x egal zero est asymptote verticale : f de x tend vers plus ou moins l infini quand x tend vers zero.} d'\'equation $x=0$ ;
			\item une \pop{asymptote horizontale}{L'axe des abscisses y egal zero est asymptote horizontale en plus et moins l infini.} d'\'equation $y=0$.
		\end{itemize}
	\end{parcoeur}
	
	\section{Application \'economique : co\^ut moyen et prix unitaire}
	
	\begin{parcoeur}[title={\ding{72}~\textbf{PAR C\OE UR}~\textbullet~Co\^ut moyen}]
		Pour $q>0$ unit\'es produites, le \textbf{co\^ut total} $C(q)$ se d\'ecompose en charges fixes $C_f>0$ et charges variables $C_v(q)$. Le \pop{co\^ut moyen}{Le cout moyen est le cout total divise par la quantite produite. Il represente le cout unitaire.} est
		\[
		CM(q)=\frac{C(q)}{q}=\frac{C_f}{q}+\frac{C_v(q)}{q}.
		\]
	\end{parcoeur}
	
	\begin{acomprendre}[title={\ding{110}~\textbf{\`A COMPRENDRE}~\textbullet~Interpr\'etation \'economique}]
		Le terme $C_f/q$ est exactement une fonction inverse : il \textbf{d\'ecro\^it} quand $q$ augmente. Les charges fixes sont r\'eparties sur un plus grand nombre d'unit\'es, donc le co\^ut unitaire baisse. C'est l'\textbf{\'economie d'\'echelle}.
	\end{acomprendre}
	
	\begin{exemplebox}[title={\ding{52}~\textbf{EXEMPLE}~\textbullet~Mod\`ele lin\'eaire + charges fixes}]
		Si $C_v(q)=c\,q$ avec $c>0$, alors $CM(q)=\dfrac{C_f}{q}+c$ : toujours d\'ecroissante, d'asymptote horizontale $y=c$.
	\end{exemplebox}
	
	\begin{exemplebox}[title={\ding{52}~\textbf{EXEMPLE}~\textbullet~Rendements d\'ecroissants}]
		Si $C_v(q)=c\,q+\dfrac{d}{q}$, alors $CM(q)=\dfrac{C_f+d}{q}+c$ : d\'ecroissante, d'asymptote horizontale $y=c$.
	\end{exemplebox}
	
	\section{Combinaison avec les polyn\^omes}
	
	\begin{acomprendre}[title={\ding{110}~\textbf{\`A COMPRENDRE}~\textbullet~Fonctions de la forme $ax+b+c/x$}]
		On consid\`ere $f(x)=ax+b+\dfrac{c}{x}$, $a\neq 0$, $c\neq 0$, d\'efinie sur $\mathbb{R}^*$ : c'est la \pop{somme d'une fonction affine et d'une fonction inverse}{Cette fonction combine une partie affine et une partie inverse. Elle est le prototype des fonctions homographiques generalisees, frequentes en economie.}.
	\end{acomprendre}
	
	\begin{parcoeur}[title={\ding{72}~\textbf{PAR C\OE UR}~\textbullet~D\'eriv\'ee et asymptote oblique}]
		\[
		f'(x)=a-\frac{c}{x^{2}}=\frac{ax^{2}-c}{x^{2}}.
		\]
		Le signe de $f'(x)$ est celui du num\'erateur $ax^{2}-c$ (le d\'enominateur $x^{2}$ est $>0$).
		\begin{itemize}[leftmargin=*,itemsep=1pt]
			\item Si $ac\le 0$ : pas de racine r\'eelle, $f'$ garde un signe constant.
			\item Si $ac>0$ : deux racines $x_1=-\sqrt{c/a}$ et $x_2=+\sqrt{c/a}$.
		\end{itemize}
		De plus $\lim_{x\to\pm\infty}\bigl[f(x)-(ax+b)\bigr]=0$ : la droite $y=ax+b$ est \pop{asymptote oblique}{Quand x tend vers plus ou moins l infini, la courbe se rapproche de la droite y egal ax plus b. Cette droite est asymptote oblique.} \`a la courbe.
	\end{parcoeur}
	
	\begin{exemplebox}[title={\ding{52}~\textbf{EXEMPLE}~\textbullet~\'Etude compl\`ete de $f(x)=x+\dfrac{1}{x}$}]
		$f'(x)=1-\dfrac{1}{x^{2}}=\dfrac{x^{2}-1}{x^{2}}=\dfrac{(x-1)(x+1)}{x^{2}}$.
		
		\begin{center}\renewcommand{\arraystretch}{1.3}
			\begin{tabular}{c|ccccccccc}
				\toprule
				$x$ & $-\infty$ & & $-1$ & & $0$ & & $1$ & & $+\infty$\\
				\midrule
				$f'(x)$ & & $+$ & $0$ & $-$ & $\|$ & $-$ & $0$ & $+$ & \\
				\midrule
				$f$ & $-\infty$ & $\nearrow$ & $-2$ & $\searrow$ & $\|$ & $\searrow$ & $2$ & $\nearrow$ & $+\infty$\\
				\bottomrule
			\end{tabular}
		\end{center}
		Asymptote oblique $y=x$ en $\pm\infty$ ; asymptote verticale $x=0$.
	\end{exemplebox}
	
	\newpage
	% ============================================================
	% VOLET 2
	% ============================================================
	\volet{Volet 2 — Tableaux \`a plusieurs entr\'ees}{Tout le cours condens\'e en tableaux comparatifs, id\'eal pour r\'eviser ou m\'emoriser.}
	
	\subsection*{Tableau 1 — R\'ecapitulatif des notions}
	
	\begin{center}\renewcommand{\arraystretch}{1.35}
		\rowcolors{2}{blue!4}{white}
		\begin{tabularx}{\linewidth}{|>{\columncolor{blue!15}}l|X|>{\columncolor{green!12}}l|}
			\hline
			\rowcolor{blue!25}\textbf{Notion} & \textbf{Contenu} & \textbf{Niveau}\\
			\hline
			Fonction & Relation $x\mapsto f(x)$ & Coll\`ege\\
			\hline
			Domaine & Ensemble des $x$ o\`u $f$ existe & Seconde\\
			\hline
			Variation & Croissance / d\'ecroissance sur un intervalle & Seconde\\
			\hline
			D\'eriv\'ee & $f'(x)=\lim_{h\to0}\frac{f(x+h)-f(x)}{h}$ & Premi\`ere\\
			\hline
			Asymptote & Droite dont la courbe s'approche ind\'efiniment & Premi\`ere\\
			\hline
			Fonction inverse & $f(x)=1/x$, d\'efinie sur $\mathbb{R}^*$ & Seconde/Premi\`ere\\
			\hline
			D\'eriv\'ee de $1/x$ & $f'(x)=-1/x^{2}$ & Premi\`ere\\
			\hline
			Co\^ut moyen & $CM(q)=C(q)/q$ & Premi\`ere STMG / G\'en\'erale\\
			\hline
			$ax+b+c/x$ & Combinaison affine + inverse & Premi\`ere/Terminale\\
			\hline
		\end{tabularx}
	\end{center}
	
	\subsection*{Tableau 2 — Domaine et parit\'e}
	
	\begin{center}\renewcommand{\arraystretch}{1.35}
		\rowcolors{2}{blue!4}{white}
		\begin{tabular}{|l|c|c|c|}
			\hline
			\rowcolor{blue!25}\textbf{Fonction} & \textbf{Expression} & $\mathcal{D}_f$ & \textbf{Parit\'e}\\
			\hline
			Inverse & $1/x$ & $\mathbb{R}^*$ & Impaire\\
			\hline
			Inverse d\'ecal\'ee & $1/(x-a)$ & $\mathbb{R}\setminus\{a\}$ & Aucune (sauf $a=0$)\\
			\hline
			Affine + inverse & $ax+b+c/x$ & $\mathbb{R}^*$ & Aucune\\
			\hline
			Homographique & $\dfrac{ax+b}{cx+d}$ & $\mathbb{R}\setminus\{-d/c\}$ & Aucune (sauf cas)\\
			\hline
			Inverse carr\'e & $1/x^{2}$ & $\mathbb{R}^*$ & Paire\\
			\hline
		\end{tabular}
	\end{center}
	
	\subsection*{Tableau 3 — D\'eriv\'ees et signes}
	
	\begin{center}\renewcommand{\arraystretch}{1.35}
		\rowcolors{2}{green!4}{white}
		\begin{tabularx}{\linewidth}{|l|c|X|c|}
			\hline
			\rowcolor{green!22}\textbf{Fonction $f$} & \textbf{D\'eriv\'ee $f'$} & \textbf{Signe de $f'$} & \textbf{Variation}\\
			\hline
			$1/x$ & $-1/x^{2}$ & Toujours $<0$ sur $\mathbb{R}^*$ & $\searrow$ sur chaque int.\\
			\hline
			$-1/x$ & $1/x^{2}$ & Toujours $>0$ sur $\mathbb{R}^*$ & $\nearrow$ sur chaque int.\\
			\hline
			$x+1/x$ & $1-1/x^{2}$ & $-$ sur $]-1,0[\cup]0,1[$, $+$ ailleurs & $\nearrow\searrow|\searrow\nearrow$\\
			\hline
			$2x+3+5/x$ & $2-5/x^{2}$ & $-$ entre $-\sqrt{5/2}$ et $\sqrt{5/2}$ & $\nearrow\searrow|\searrow\nearrow$\\
			\hline
			$ax+b+c/x$ & $a-c/x^{2}$ & Signe de $ax^{2}-c$ & D\'epend du signe de $ac$\\
			\hline
		\end{tabularx}
	\end{center}
	
	\subsection*{Tableau 4 — Limites et asymptotes}
	
	\begin{center}\renewcommand{\arraystretch}{1.35}
		\rowcolors{2}{orange!5}{white}
		\begin{tabular}{|l|c|l|}
			\hline
			\rowcolor{orange!25}\textbf{Limite} & \textbf{R\'esultat} & \textbf{Cons\'equence g\'eom\'etrique}\\
			\hline
			$\lim_{x\to+\infty} 1/x$ & $0^{+}$ & Asymptote horizontale $y=0$ (c\^ot\'e $+$)\\
			\hline
			$\lim_{x\to-\infty} 1/x$ & $0^{-}$ & Asymptote horizontale $y=0$ (c\^ot\'e $-$)\\
			\hline
			$\lim_{x\to 0^{+}} 1/x$ & $+\infty$ & Asymptote verticale $x=0$ (c\^ot\'e $+$)\\
			\hline
			$\lim_{x\to 0^{-}} 1/x$ & $-\infty$ & Asymptote verticale $x=0$ (c\^ot\'e $-$)\\
			\hline
			$\lim_{x\to\pm\infty}[f(x)-(ax+b)]$ & $0$ & Asymptote oblique $y=ax+b$\\
			\hline
		\end{tabular}
	\end{center}
	
	\subsection*{Tableau 5 — Application \'economique}
	
	\noindent Avec $C_f=1000$~EUR et $c=2$~EUR/unit\'e : $CM(q)=1000/q+2$.
	
	\begin{center}\renewcommand{\arraystretch}{1.3}
		\rowcolors{2}{blue!4}{white}
		\begin{tabular}{|r|c|c|c|c|c|c|}
			\hline
			\rowcolor{blue!22}$q$ (unit\'es) & $1$ & $5$ & $10$ & $50$ & $100$ & $1000$\\
			\hline
			$C_f/q$ (EUR) & $1000$ & $200$ & $100$ & $20$ & $10$ & $1$\\
			\hline
			$CM(q)$ (EUR) & $1002$ & $202$ & $102$ & $22$ & $12$ & $3$\\
			\hline
		\end{tabular}
	\end{center}
	
	\noindent\textbf{Lecture.} $q$ augmente $\Rightarrow C_f/q\to 0$ et $CM(q)\to c=2$ : c'est l'asymptote horizontale $y=2$.
	
	\subsection*{Tableau 6 — Comparaison de fonctions usuelles}
	
	\begin{center}\renewcommand{\arraystretch}{1.3}
		\rowcolors{2}{purple!4}{white}
		\begin{tabular}{|l|c|c|c|c|}
			\hline
			\rowcolor{purple!22}\textbf{Fonction} & \textbf{Domaine} & \textbf{D\'eriv\'ee} & \textbf{Variation} & \textbf{Asymptotes}\\
			\hline
			$x$ & $\mathbb{R}$ & $1$ & $\nearrow$ & aucune\\
			\hline
			$x^{2}$ & $\mathbb{R}$ & $2x$ & $\searrow$ puis $\nearrow$ & aucune\\
			\hline
			$1/x$ & $\mathbb{R}^*$ & $-1/x^{2}$ & $\searrow$ sur chaque int. & $x=0$, $y=0$\\
			\hline
			$\sqrt{x}$ & $[0,+\infty[$ & $1/(2\sqrt{x})$ & $\nearrow$ & aucune\\
			\hline
			$x+1/x$ & $\mathbb{R}^*$ & $1-1/x^{2}$ & $\nearrow\searrow|\searrow\nearrow$ & $x=0$, $y=x$\\
			\hline
		\end{tabular}
	\end{center}
	
	\newpage
	% ============================================================
	% VOLET 3 — CARTE MENTALE (portrait)
	% ============================================================
	\volet{Volet 3 — Carte mentale interactive}{Survolez chaque n\oe ud pour afficher son contenu d\'etaill\'e.}
	
	\begin{center}
		\begin{tikzpicture}[
			scale=0.85, transform shape,
			central/.style={circle,draw=blue!70!black,fill=blue!25,thick,
				minimum size=2.8cm,align=center,font=\bfseries\footnotesize,inner sep=2pt},
			branch/.style={rounded corners=4pt,draw=green!50!black,fill=green!20,thick,
				minimum width=2.4cm,minimum height=0.95cm,align=center,
				font=\footnotesize\bfseries,inner sep=3pt},
			leaf/.style={rounded corners=3pt,draw=orange!70!black,fill=orange!15,
				align=center,font=\scriptsize,inner sep=2pt,minimum height=0.6cm},
			edge/.style={thick,blue!50!black},
			ledge/.style={orange!70!black,thick}
			]
			\node[central] (C) at (0,0)
			{\tip{\shortstack{Fonction\\inverse\\$x\mapsto 1/x$}}%
				{Fonction f de x egale 1 sur x, definie sur R prive de 0. Impaire. Decroissante sur chaque intervalle. Derivee moins 1 sur x carre. Asymptotes x egal 0 et y egal 0.}};
			
			\node[branch] (Bdef)  at (90:6.0)
			{\tip{D\'efinition}{f de x egale 1 sur x, definie sur R prive de 0}};
			\node[branch] (Bder)  at (30:6.0)
			{\tip{D\'eriv\'ee}{f prime de x egale moins 1 sur x carre}};
			\node[branch] (Bvar)  at (-30:6.0)
			{\tip{Variations}{Decroissante sur moins infini zero et sur zero plus infini}};
			\node[branch] (Basy)  at (-90:6.0)
			{\tip{Asymptotes}{x egal zero : verticale. y egal zero : horizontale}};
			\node[branch] (Beco)  at (-150:6.0)
			{\tip{\'Economie}{Cout moyen, prix unitaire}};
			\node[branch] (Bpoly) at (150:6.0)
			{\tip{Polyn\^omes}{Fonctions de la forme ax plus b plus c sur x}};
			
			\foreach \b in {Bdef,Bder,Bvar,Basy,Beco,Bpoly}
			\draw[edge] (C) -- (\b);
			
			\node[leaf] (Ld1) at (105:10)  {\tip{$f(x)=1/x$}{Pour tout x non nul, f de x egale 1 sur x}};
			\node[leaf] (Ld2) at (90:10.5) {\tip{$\mathcal{D}_f=\mathbb{R}^*$}{R prive de zero : reunion de moins infini zero et zero plus infini}};
			\node[leaf] (Ld3) at (75:10)   {\tip{Impaire}{f de moins x egale moins f de x : symetrie centrale par rapport a l origine}};
			\foreach \l in {Ld1,Ld2,Ld3} \draw[ledge] (Bdef) -- (\l);
			
			\node[leaf] (Lr1) at (45:10)   {\tip{$f'(x)=-1/x^{2}$}{Derivee de x puissance moins 1}};
			\node[leaf] (Lr2) at (30:10.5) {\tip{Signe $<0$}{Toujours negative sur R prive de 0}};
			\node[leaf] (Lr3) at (15:10)   {\tip{$(1/v)'=-v'/v^{2}$}{Formule generale de derivation d un quotient}};
			\foreach \l in {Lr1,Lr2,Lr3} \draw[ledge] (Bder) -- (\l);
			
			\node[leaf] (Lv1) at (-15:10)  {\tip{$\searrow$ sur $]-\infty,0[$}{f prime strictement negative}};
			\node[leaf] (Lv2) at (-30:10.5){\tip{$\searrow$ sur $]0,+\infty[$}{f prime strictement negative}};
			\node[leaf] (Lv3) at (-45:10)  {\tip{Pas globalement}{Contre-exemple : moins 1 est inferieur a 1 mais f de moins 1 est inferieur a f de 1}};
			\foreach \l in {Lv1,Lv2,Lv3} \draw[ledge] (Bvar) -- (\l);
			
			\node[leaf] (La1) at (-75:10)  {\tip{$x=0$ verticale}{limite plus ou moins l infini quand x tend vers zero}};
			\node[leaf] (La2) at (-90:10.5){\tip{$y=0$ horizontale}{limite zero quand x tend vers plus ou moins l infini}};
			\node[leaf] (La3) at (-105:10) {\tip{Pas d'oblique}{Sauf pour ax plus b plus c sur x : asymptote y egal ax plus b}};
			\foreach \l in {La1,La2,La3} \draw[ledge] (Basy) -- (\l);
			
			\node[leaf] (Le1) at (-135:10) {\tip{Co\^ut moyen}{CM de q egale C f sur q plus c}};
			\node[leaf] (Le2) at (-150:10.5){\tip{Prix unitaire}{Terme en 1 sur q frequent en economie}};
			\node[leaf] (Le3) at (-165:10) {\tip{\'Econ.\ d'\'echelle}{CM baisse quand q augmente}};
			\foreach \l in {Le1,Le2,Le3} \draw[ledge] (Beco) -- (\l);
			
			\node[leaf] (Lp1) at (165:10)  {\tip{$ax+b+c/x$}{Combinaison affine plus inverse}};
			\node[leaf] (Lp2) at (150:10.5){\tip{$f'=a-c/x^{2}$}{Numerateur ax carre moins c, denominateur x carre}};
			\node[leaf] (Lp3) at (135:10)  {\tip{Asymptote oblique}{y egal ax plus b en plus ou moins l infini}};
			\foreach \l in {Lp1,Lp2,Lp3} \draw[ledge] (Bpoly) -- (\l);
			
		\end{tikzpicture}
	\end{center}
	
	\vfill
	\begin{acomprendre}[title={\ding{110}~\textbf{\`A COMPRENDRE}~\textbullet~Lecture de la carte}]
		La carte regroupe les \textbf{six axes essentiels} de la le\c{c}on : \emph{d\'efinition}, \emph{d\'eriv\'ee}, \emph{variations}, \emph{asymptotes}, \emph{\'economie}, \emph{polyn\^omes}. Pour chaque axe, trois propri\'et\'es-cl\'es sont d\'evelopp\'ees. Survolez un n\oe ud pour obtenir son contenu d\'etaill\'e.
	\end{acomprendre}
	
	\newpage
	% ============================================================
	% VOLET 4 — GRAPHE D'INTERACTIONS (portrait)
	% ============================================================
	\volet{Volet 4 — Graphe des interactions}{Survolez chaque n\oe ud pour afficher son r\^ole d\'etaill\'e dans la le\c{c}on.}
	
	\begin{center}
		\begin{tikzpicture}[
			scale=0.85, transform shape,
			concept/.style={rounded corners=6pt,draw=blue!70!black,fill=blue!20,thick,
				minimum width=2.6cm,minimum height=1cm,align=center,font=\footnotesize\bfseries,
				inner sep=3pt},
			app/.style={rounded corners=6pt,draw=orange!70!black,fill=orange!20,thick,
				minimum width=2.6cm,minimum height=1cm,align=center,font=\footnotesize\bfseries,
				inner sep=3pt},
			mathrel/.style={rounded corners=6pt,draw=green!50!black,fill=green!20,thick,
				minimum width=2.6cm,minimum height=1cm,align=center,font=\footnotesize\bfseries,
				inner sep=3pt},
			arrow/.style={-{Latex[length=2.5mm]},thick,blue!60!black},
			arrowapp/.style={-{Latex[length=2.5mm]},thick,orange!70!black,dashed},
			arrowmath/.style={-{Latex[length=2.5mm]},thick,green!50!black},
			lbl/.style={font=\tiny,fill=white,inner sep=1pt,rounded corners=1pt}
			]
			\node[concept] (FI)  at (0,0)    {\tip{Fonction inverse}{f de x egale 1 sur x, definie sur R prive de 0, impaire}};
			\node[mathrel] (DEF) at (0,5)    {\tip{D\'efinition / domaine}{Domaine R prive de 0. Impaire. Courbe hyperbole}};
			\node[mathrel] (DER) at (7,3.3)  {\tip{D\'eriv\'ee}{f prime egale moins 1 sur x carre, strictement negative}};
			\node[mathrel] (VAR) at (8.5,0)  {\tip{Variations}{Decroissante sur chaque intervalle du domaine}};
			\node[mathrel] (ASY) at (7,-3.3) {\tip{Asymptotes}{Verticale x egal zero, horizontale y egal zero}};
			\node[concept] (COU) at (0,-5)   {\tip{Hyperbole}{Courbe de la fonction inverse, symetrique par rapport a l origine}};
			\node[app] (ECO) at (-7,-3.3)    {\tip{\'Economie}{Cout moyen, prix unitaire : CM de q egale C f sur q plus c}};
			\node[app] (POL) at (-7,3.3)     {\tip{Polyn\^omes}{Fonctions ax plus b plus c sur x}};
			
			\draw[arrowmath] (DEF) -- node[lbl]{permet} (DER);
			\draw[arrowmath] (DER) -- node[lbl,above right]{signe} (VAR);
			\draw[arrowmath] (VAR) to[bend left=8] node[lbl,above]{allure} (COU);
			\draw[arrowmath] (ASY) to[bend right=8] node[lbl,below]{branches} (COU);
			\draw[arrowmath] (DEF) to[bend right=15] node[lbl,left]{limites} (ASY);
			\draw[arrow]     (FI)  -- node[lbl,below left]{mod\'elise} (ECO);
			\draw[arrow]     (FI)  -- node[lbl,above left]{se combine} (POL);
			\draw[arrowapp]  (POL) -- node[lbl,left]{d\'erivation} (DER);
			\draw[arrowapp]  (ECO) -- node[lbl,below]{\'etude} (VAR);
			\draw[arrowapp]  (POL) to[bend left=20] node[lbl,left]{oblique} (ASY);
			\draw[arrowmath] (DEF) to[bend left=20] node[lbl,above]{applique} (FI);
			
			\node[font=\scriptsize,align=left,anchor=north west] at (-10.5,-6.8) {%
				\tikz\draw[arrowmath] (0,0)--(0.6,0);\ math\'ematiques\\
				\tikz\draw[arrow] (0,0)--(0.6,0);\ lien direct\\
				\tikz\draw[arrowapp] (0,0)--(0.6,0);\ application};
		\end{tikzpicture}
	\end{center}
	
	\vspace{0.4cm}
	\begin{acomprendre}[title={\ding{110}~\textbf{\`A COMPRENDRE}~\textbullet~Lecture du graphe}]
		Le graphe met en \'evidence les \textbf{interactions} entre les notions. Par exemple : la \emph{d\'efinition} permet le calcul de la \emph{d\'eriv\'ee}, dont le signe donne les \emph{variations}, qui elles-m\^emes d\'eterminent l'\emph{allure de la courbe} ; r\'eciproquement, les \emph{limites} issues de la d\'efinition donnent les \emph{asymptotes}, qui sont des propri\'et\'es de la courbe. Les arcs \textcolor{orange!80!black}{orange} relient la fonction inverse \`a ses \textbf{applications} (\'economie, combinaison polynomiale), et montrent comment ces applications r\'einvestissent l'arsenal math\'ematique (d\'eriv\'ee, variations, asymptotes).
	\end{acomprendre}
	
\end{document}